\section{Method: how speed was measured and correctness enforced}\label{sec:method}

\subsection{Models}
Six models span the sizes and features that matter (Table~\ref{tab:models}). Four are the openEMS Python tutorials,
two of them with the mesh refined (``x5'', ``x3'') to a size that is clearly bandwidth-bound; two are the author's
real boards, built by the author's toolkit. Every feature in daily use is covered: Mur and PML boundaries (full-z slabs
and a z-only slab), lumped, microstrip-line (MSL) and coaxial ports, conducting sheets (the dispersive-material
extension), near-field-to-far-field (nf2ff) dumps, frequency-domain field dumps. The cylindrical multigrid was checked
separately on openEMS's CylindricalWave\_CC tutorial.

\begin{table}[H]\centering\small
\resizebox{\textwidth}{!}{\begin{tabular}{@{}llrl@{}}
\toprule
model & origin & cells & features exercised \\
\midrule
patch tutorial & Simple\_Patch\_Antenna.py & 0.104\,M & lumped port, 6$\times$ MUR, nf2ff box (12 time-domain dumps) \\
notch tutorial & MSL\_NotchFilter.py & 0.270\,M & MSL ports, PML\_8 at both x ends, PEC at $z_{\min}$, MUR elsewhere \\
patch x5 & patch, mesh $\times$5 & 7.86\,M & bandwidth-bound size, MUR \\
notch x3 & notch, mesh $\times$3 & 5.25\,M & bandwidth-bound size, PML slabs \\
GCPW line (kgcpw) & rfsim gcpw\_line & 0.242\,M & conducting sheets (dispersive extension), MSL ports, MUR, 17 z lines \\
SMA launch (ksma) & rfsim sma\_launch & 1.071 / 0.871\,M & coaxial port, z-only PML\_8, 69\,500 / 64\,410 distinct operators, FD dumps \\
\bottomrule
\end{tabular}}
\caption{Benchmark models. ``M'' = million Yee cells. The SMA launch has two meshes: 1.071\,M cells with the
toolkit's former 0.02\,mm mesh floor, the mesh of the benchmark tables up to snapshot cg18, and 0.871\,M with the
0.05\,mm floor that is the toolkit's default since 2026-11-04 (toolkit commit 4a7e20e, Section~\ref{sec:meshfloor}),
on which the benchmark model has been built since then: the GPU engine's
measurements (Section~\ref{sec:gpu}), the other engines' runs (Section~\ref{sec:positioning}), the checks of cg29 and
its comparison with cg18 (\code{results/cpu\_items\_cg29\_2026-11-05.txt}) use it. Tables that use the 1.071\,M mesh
say so. Rates on the two meshes do not compare directly. Operator counts: 69\,500 on the 1.071\,M mesh (project log,
not in \code{results/}), 64\,410 on the 0.871\,M mesh
(\code{results/diamond\_engine\_cg24\_cg27\_2026-11-04.txt}).}\label{tab:models}
\end{table}

\subsection{The harness}
\code{bench.py MODEL --configs CFG...} builds the model in a fresh Python process per run (\code{runner.py}) and calls
\code{FDTD.Run()} with a fixed number of time steps (\code{EndCriteria=0}), so every engine performs identical work. A
configuration names the library to load (\code{SNAP@}, a saved set of \code{libopenEMS}, \code{libCSXCAD},
\code{libnf2ff} dylibs picked up through \code{DYLD\_LIBRARY\_PATH}), the engine, the thread count, and environment
variables for individual knobs (\code{\#NAME=VAL}). The first configuration is the reference; for every other run the
harness reports the maximum relative difference of the port probe files and the largest difference of
$S_{11}$/$S_{21}$ (linear and dB, over the frequencies where $|S_{\rm ref}|>-40$\,dB), the engine speed in \MC\ parsed
from openEMS's own output, and the CPU load of \emph{other} processes in the second before the run (\code{others NN\%},
from \code{top}). A probe file whose reference is all zero, and an $S$-parameter that stays below $-40$\,dB, cannot be
compared that way, and the published tables leave them out without a note: the notch x3 check ends before the pulse
reaches port 2, so its tables compare neither the port-2 probes nor $S_{21}$, and the GCPW tables do not compare
$S_{11}$. \code{--quiet 60} waits until other processes use less than 60\,\% of a core, for at most \code{--quiet-max}
seconds; \code{--repeat 2} runs each configuration twice and reports the faster run (both speeds are listed), and in
the tables in \code{results/} only each configuration's last run was compared. The harness now compares every probe
file and $S$-parameter exactly, over all frequencies, and names those without a figure in its output; it also compares
every run with the reference configuration's first run and flags the reported run when its quiet wait timed out.
Since 2026-11-04 the two GPU engines (\code{gpu}, \code{metal-diamond}) and any configuration with
\code{OPENEMS\_FTZ} set count as tolerant: their differences from a stock reference are printed but do not fail the
run, since neither is bit-identical to a run that keeps subnormal floats (the GPU engines' reference is the CPU
engine run with \code{OPENEMS\_FTZ=1}; Section~\ref{sec:policy}). From the cg27 tables of the same day on, each line
also gives the CPU and GPU use of the other processes during the run (\code{bench/usage.py}), following the author's
rule of that day that every benchmark line carries both. A knob's value cannot contain a comma, because the harness
splits its settings on commas; a tile setting \code{M,V} of the GPU engine was once silently dropped that way (the
engine now also reads \code{M:V}). From cg24 (2026-11-04) most speed comparisons are paired and interleaved (\code{bench/abtest.py}):
the configurations run in turn, round after round, and a ratio is the median over the rounds of the per-round ratio
of two configurations, given with its range, so that a slow phase of the shared machine hits both configurations of
a round alike. The unpaired runs of 2026-11-04 and 2026-11-05 (the faster of two runs, or one run each: the tile
sweep of cg23, the other engines, the Swift diamond mode, part B of the $K=1$ file, the GPU end-to-end runs of cg23
and the check tables) say so where they are quoted.
Runs go through
\code{FDTD.Run()} rather than \code{Write2XML} plus the binary because in v0.37.0 (with CSXCAD v0.7.0), and so in this
fork, \code{Write2XML} writes primitive coordinates with 7 significant digits but mesh lines with 15, so
zero-thickness boxes (an MSL port's excitation plane) fall off their mesh line and are dropped, reported only as
unused primitives. Upstream master has since added a drawing tolerance taken from the mesh (commit db68e62,
2026-10-21) that covers this case.

\code{check.sh TAG SNAP} is the acceptance test: the six models at 1, 4, 8 and 10 threads and with the automatic
thread count, against the stock snapshot, two repeats, results in the top-level \code{results/} folder as
\code{MODEL\_sS\_nN\_TAG.txt}. A snapshot was accepted only when every comparison printed \code{probes rel=0.0e+00} and
\code{max|dS|=0.0e+00}; \code{bench.py} and \code{check.sh} now also exit with an error status when a comparison
differs or a file is missing. cg18 (patches 0001--0006, the build installed on the author's Mac from 2026-11-04 to
2026-11-06) passed the full check on a quiet machine on 2026-11-03 (\code{results/*\_cg18.txt},
Section~\ref{sec:results}). The current build, cg29 (patches 0001--0027, installed in cg18's place on 2026-11-06),
passed it on 2026-11-05 with three more configurations, the geometry coefficient tables
forced at 4 threads, the packed ones at 8 and the barriers instead of the neighbour flags at 8
(\code{results/*\_cg29.txt}; the speeds in those tables are the faster of two unpaired runs under the check's load,
not comparisons).
Its other checks (\code{results/cpu\_items\_cg29\_2026-11-05.txt}, section 3): \code{scripts/checks/debye\_check.sh};
the validation suite, 15 of 15 with the multithreaded engine at 4 threads and with the fastest engine; the cylindrical
multigrid, 34 datasets equal to stock's; box models with conducting sheets, Drude/Lorentz blocks, UPML, Mur faces and
the steady-state detection against the basic or sse-compressed engine of the same build, whose extension code is
stock's; and a matrix of 45 small models in up to 11 run modes, 461 comparisons identical and 10 that differ only in
the stop step of one model, for the reason given under ``Energy end criterion'' below. The Metal engine of
Section~\ref{sec:gpu} has its own acceptance test, against the multithreaded engine run with \code{OPENEMS\_FTZ=1}
instead of stock (Section~\ref{sec:policy}): \code{scripts/checks/diamond\_check.sh} compares every exported field
value of box models (PEC, Mur and PML faces, delayed and soft excitations, lumped RLC elements, conducting sheets,
Drude/Lorentz and Debye blocks) at $K=1$, 2, 3, 4 and 8 time steps per tile and with an odd tile, threadgroup and
batch size, 72 of 72 comparisons identical on cg28 and cg29, and \code{bench.py} compares the port probes and
$S$-parameters of the six models, identical on cg29 (\code{results/*\_cg29gpu.txt}) as on the earlier GPU snapshots.
The cg29 runs set \code{OPENEMS\_FTZ=1} for the whole call (\code{results/cpu\_items\_cg29\_2026-11-05.txt},
section 3), so the reference is tolerant too, and a tolerant reference makes every configuration tolerant: the files
end with \code{NOTHING REQUIRED TO MATCH}, and the identity is in each line's \code{probes rel=0.0e+00} and
\code{max|dS|=0.0e+00}.
When cg29 was installed on 2026-11-06 (its openEMS library and \code{openems.h} replaced cg18's; the install's CSXCAD
and nf2ff libraries already had cg29's bytes; the Python modules stayed the v0.37.0 binding, which affects SAR but
not FDTD runs or nf2ff, \code{BUILD.md} section 6), it was checked from the install with the command line of
Section~\ref{sec:repro}: the reduced check (\code{check.sh} with \code{LITE=1}: 4 and 8 threads and the automatic
count, one run each) identical to stock on all six models, \code{debye\_check.sh} passed (with the same fields as the
cg29 snapshot), \code{diamond\_check.sh} 72 of 72, the validation suite 15 of 15 with the fastest engine and with the
multithreaded engine at 4 threads, and the GCPW end-to-end run with the toolkit's defaults stopped at step 34\,300, as
on cg18 (Table~\ref{tab:e2e}). Other processes used 4.6--8.5 cores during these runs, so their speeds are not quoted,
and their tables are not in \code{results/}; \code{Docs/fdtdKit/PLAN.md}, section 16, logs the runs.
On 2026-11-07 the 27 published patches were built from a fresh clone of upstream's superproject into an empty prefix
(\code{BUILD.md}, sections 1, 4 and 5; snapshot pub29), and that build passed \code{check.sh} (every run identical
to stock on all six models, \code{results/*\_pub29.txt}), \code{debye\_check.sh}, \code{diamond\_check.sh} 72 of
72, the six models on the Metal engine, the validation suite 15 of 15 with the multithreaded engine at 4 threads and
with the fastest engine, and the cylindrical multigrid (\code{results/fresh\_build\_pub29\_2026-11-07.txt}); other
programs used up to 8.9 cores during \code{check.sh}, so here too only the verdicts count.

\subsection{Verification policy}\label{sec:policy}
\begin{itemize}
\item \textbf{Bit identity, not tolerance.} Floating-point FDTD is sensitive enough to rounding that a reordered sum
  shows up at $10^{-7}$ after a few thousand steps and grows. Every CPU change keeps the \emph{same} float operations
  in the \emph{same} order per cell; only who executes them, when, and from which copy of the coefficients changes.
  Hence the comparison is exact, and it caught a real bug: a hook plan built before the cylindrical multigrid
  extension registered, which made two snapshots wrong on that example only. The Metal engine of
  Section~\ref{sec:gpu} is held to the same rule where the arithmetic can match: its kernels give every cell the CPU
  engines' operations in the same order, with \code{fma()} exactly where clang fuses on arm64 and Metal's fast math
  off, so its fields must equal those of the multithreaded engine run with \code{OPENEMS\_FTZ=1}, which flushes
  subnormal floats to zero as Apple's GPU does. Where the arithmetic cannot match (a GPU engine against a CPU run that
  keeps subnormals, stock included, or another implementation), the harness reports the largest relative difference
  of the port probes and the largest $S$-parameter difference in dB, and these are measurements, not acceptance
  criteria (``Named exceptions'' below).
\item \textbf{Field dumps too.} HDF5 time- and frequency-domain dumps are compared dataset by dataset with
  \code{scripts/checks/h5cmp.py}, run by hand rather than by the harness. On the cylindrical example the earlier
  checks compared 18 datasets (project log); the check of cg18, a run of 1\,500 steps with a frequency- and a
  time-domain dump, compared 34, all equal to stock's (\code{results/correctness\_cg18\_2026-11-03.txt}), and so did
  that of cg29 (\code{results/cpu\_items\_cg29\_2026-11-05.txt}, section 3).
\item \textbf{Named exceptions.} No GPU engine measured here is bit-identical to stock. Upstream's GPU engine in the
  fork's base (\code{--engine=gpu}, the pre-merge \code{gpu-engine} branch, commit 7bb294f) differs from the CPU
  engines by up to $2.3\times10^{-4}$\,dB in $S_{11}$ (notch tutorial, \code{results/notch\_s1\_n16000\_cg2.txt}).
  The cause is not the order of operations, which the Metal kernels follow operation by operation, but that Apple's
  GPU flushes subnormal floats to zero while the CPU cores keep them (Section~\ref{sec:beyond}): in a 3\,000-step
  check on the notch tutorial (snapshot cg17, run in the final review of 2026-11-03,
  \code{results/notch\_s1\_n3000\_ftzgpu.txt}), the multithreaded engine at 2 threads and the sse-compressed engine,
  both with \code{OPENEMS\_FTZ=1}, gave port probes and $S_{11}$ identical to the GPU engine's, while the
  multithreaded engine without it differed by $3.2\times10^{-7}$ relative in the probes and $9.3\times10^{-6}$\,dB in
  $S_{11}$; field dumps were not compared in that check. On 2026-11-04 the same held on all six models: with
  \code{OPENEMS\_FTZ=1} on the CPU, the base's GPU engine and the Metal diamond engine both gave port probes and
  $S$-parameters identical to the multithreaded engine's (\code{results/*\_cg22gpu.txt},
  \code{results/diamond\_engine\_cg22\_cg23\_2026-11-04.txt}). Against stock without the flush the diamond engine
  differs by $1.7\times10^{-7}$ to $9.4\times10^{-7}$ relative in the port probes, by at most
  $2.3\times10^{-4}$\,dB in $S_{11}$ (notch tutorial; at most $2.3\times10^{-5}$\,dB on the other models) and by at
  most $9.4\times10^{-5}$\,dB in $S_{21}$ (\code{results/*\_cg22.txt}; its results have not changed since). Its nf2ff
  time-domain dumps were compared too: identical on the patch tutorial (984 datasets) and patch x5 (120) when both
  runs set \code{OPENEMS\_FTZ=1}, because the dumps are computed on the creating thread, which the switch also
  flushes; with the flush in the CPU run only, 27 of the tutorial's 984 datasets differed while the fields were
  identical (\code{results/diamond\_engine\_cg20\_2026-11-04.txt}). Upstream master's reworked GPU engine (merged as
  c3a956b on 2026-10-17) was measured on 2026-11-04 in a build of master 5fecbb5 (\code{results/others\_2026-11-04/},
  Section~\ref{sec:positioning}): on five of the six models its differences from stock are, to the printed digits,
  those of the diamond engine above; on the SMA launch master's CPU engine differs from stock as well
  ($1.5\times10^{-4}$ relative in the probes, $3.4\times10^{-3}$\,dB in $S_{11}$), a change in master since v0.37.0
  that its GPU engine shows too. It was not compared with a CPU run with \code{OPENEMS\_FTZ=1}. The new nf2ff path
  accumulates in BLAS order and differs from the legacy loop by up to $7\times10^{-7}$ of the pattern maximum in the
  project log, and by $4.5$--$6.8\times10^{-7}$ on the patch tutorial with cg18
  (\code{results/correctness\_cg18\_2026-11-03.txt}); a weak cross-polar component can differ more relative to its own
  peak. The legacy loop is kept behind \code{OPENEMS\_NF2FF\_LEGACY=1} and is itself bit-identical (also checked in
  that file).
\item \textbf{Energy end criterion.} The fork's base checks the energy end criterion every Nyquist period during the
  first minute of a run (upstream commit eb1c5ee, in the \code{gpu-engine} base, not in v0.37.0); v0.37.0 checks it,
  unless run with \code{--exact-endcriteria}, only at its progress report every 4\,s of wall time, so its stop step
  depends on the machine's speed. A run that ends on the energy criterion can therefore stop at a different step than
  stock v0.37.0. For the same number of time steps the results are identical, which is why every comparison here uses
  a fixed step count. In the run-mode matrix of the cg29 check this was the only difference: on a 0.7\,s model that
  ends on the energy criterion, stock ran to its step limit of 1\,496 steps and every cg29 configuration stopped at
  step 255, with its 64 probe lines equal to the first 64 of stock's 375
  (\code{results/cpu\_items\_cg29\_2026-11-05.txt}, section 3).
\item \textbf{Every number has a load figure next to it.} The machine was shared with other sessions; a speed quoted
  without its \code{others} percentage was not trusted, and comparisons across snapshots were re-run when loads
  differed. Up to the cg23 tables of 2026-11-04 the figure was a one-second sample taken just before each run, not
  during it, and it failed once: the snapshot-A row at 8 threads in \code{results/patch\_s5\_n1200\_AB.txt} reads
  \code{others -15493379\%}, so the load during that run is unknown. From the cg27 tables of 2026-11-04 on
  (\code{bench/usage.py}) the lines also give the other processes' CPU and GPU use during the run, and from cg24 on
  most speed comparisons are paired and interleaved (the harness, above); the unpaired ones say so where they are
  quoted.
\item \textbf{Nothing is accepted on speed alone.} Several items can be switched off by an environment variable that
  restores the previous path (listed in \code{BUILD.md}): \code{OPENEMS\_FUSED=0} (fused sweep),
  \code{OPENEMS\_BALANCE=0} (slab balancing), \code{OPENEMS\_QOS=0} (thread QoS), \code{OPENEMS\_NO\_STAGGER}
  (allocator), \code{OPENEMS\_NO\_PML\_TABLE} (PML tables), \code{OPENEMS\_NF2FF\_LEGACY} (far-field product) and
  \code{OPENEMS\_NO\_PACKED\_COEFF}, which turns off three items at once (Section~\ref{sec:engine} says which). cg29
  adds \code{OPENEMS\_OP\_INDEX\_PER\_CELL} (the former operator index per cell), \code{OPENEMS\_ADE\_MERGED=0} (the
  ADE orders of a conducting sheet one after the other), \code{OPENEMS\_PML\_YVEC=0} (the scalar UPML rows one at a
  time), \code{OPENEMS\_COEF} (\code{packed} or \code{geom}: the coefficient layout), \code{OPENEMS\_FLAGS=0} (the
  two barriers instead of the neighbour flags) and \code{OPENEMS\_FLAGS\_SPIN\_US} (the flag spin at any thread
  count; Section~\ref{sec:cg29}). The \code{OPENEMS\_NO\_*}, \code{OPENEMS\_NF2FF\_LEGACY},
  \code{OPENEMS\_OP\_INDEX\_PER\_CELL} and \code{OPENEMS\_FTZ} switches act when set to any value other than empty or
  0; \code{OPENEMS\_FTZ} and the probe end criterion are off unless set. The other items (the spin barriers and the
  hook plan, work stealing and the thread search, the Mur and UPML fast paths, the dispersive-material line hooks, the
  dump gather, the setup changes and the CSXCAD changes, and from cg29 the threaded UPML coefficient setup and the
  line hooks of the steady-state detection and the Debye materials) have no switch: they rest on the bit-identity
  checks and are bisected through the snapshots and commits. The Metal engine is opt-in
  (\code{--engine=metal-diamond}); \code{OPENEMS\_DIAMOND\_TUNE=0} turns off its run-time choice of the tiling, and
  \code{OPENEMS\_DIAMOND\_K}, \code{OPENEMS\_DIAMOND\_TILE} and \code{OPENEMS\_DIAMOND\_COEF} fix the tiling and the
  coefficient layout (Section~\ref{sec:gpu}).
\item \textbf{What ``bit-identical'' rests on.} Three rewrites depend on the compiler fusing a multiply-add the same
  way in the old and the new expression; this was verified by comparison on every model with Apple clang builds, not
  forced in the source (Section~\ref{sec:contract}). In the standalone kernel the rewritten inner loop matches the stock
  loop bit for bit under Apple clang but not under GCC~16 (\code{results/kernel\_gcc\_vs\_clang\_2026-11-03.txt}).
  The Metal kernels write \code{fma()} where Apple clang fuses the CPU expressions, so the GPU engine's identity rests
  on the same compiler's choices.
\item \textbf{The policy has a known hole:} a bit-identical test only covers the features the test models exercise;
  the Debye material of Section~\ref{sec:reviewfixes} was not in them and its bug surfaced only in the
  pre-publication review.
\end{itemize}

\subsection{Baseline: stock openEMS v0.37.0}\label{sec:baseline}
Table~\ref{tab:step1} is the first measurement. Class and engine names in this paper are those of v0.37.0 and of the
fork's base, upstream's \code{gpu-engine} branch at 7bb294f. Upstream master has since reorganised them (commit
c0816a2, 2026-11-02): \code{Operator\_SSE\_Compressed} became \code{Coeffs\_SSE\_Compressed},
\code{Engine\_SSE\_Compressed} and \code{Engine\_Multithread} gave way to \code{Engine\_sse} running on an
\code{Engine\_ThreadPool}, and \code{--engine=sse}, \code{sse-compressed} and \code{multithreaded} now all select that
engine. Three observations shaped the work:
\begin{enumerate}
\item The multithreaded engine's best thread count is 8 on the big models but 2--4 on the small ones, and the small
  ones get \emph{slower} with every thread past that. The per-time-step synchronisation is the reason. Stock runs
  $2+6N$ barriers per step for $N$ extensions: 50 on the patch model (one excitation, six Mur faces, and a lumped-RLC
  extension without active cells that the port's feed resistor creates) and 38 on the notch. The initial review
  (below; it counted 44 for the patch, missing that extension) measured 42 such barriers at 0.21\,ms per step on 4
  threads and 0.66\,ms on 8 on the tutorial grid, where a step is 0.09\,ms of work, and at 0.36 and 1.0\,ms on the
  7.9-million-cell grid.
\item The default \code{numThreads=0} starts with one thread and adds one every $\sim$4\,s of run time, so a 10--20\,s
  run spends most of its life under-threaded. The ``fastest'' column is 41\,\% and 37\,\% slower than the best fixed
  count on the two large models (15.3\,s and 19.6\,s runs), 22\,\% on the notch tutorial (17.5\,s) and 8\,\% on the
  patch tutorial (4.3\,s).
\item PML slabs at both x ends load the first and last thread's slab; the notch x3 model is flat at 671--693\,\MC\ for
  3--6 threads and only reaches 845 at 8.
\end{enumerate}

\begin{table}[H]\centering\small
\resizebox{\textwidth}{!}{\begin{tabular}{@{}lrrrrrr@{}}
\toprule
model & basic & sse & sse-compressed & multithreaded best (thr) & multithreaded 4 & default (``fastest'') \\
\midrule
patch x5 (7.86\,M)    & 157 & 265 & 394 & 1041 (8) & 1006 & 615 \\
notch x3 (5.25\,M)    & 147 & 204 & 290 &  845 (8) &  693 & 534 \\
patch tutorial (0.10\,M) & 162 & 195 & 255 & 277 (2) & 205 & 254 \\
notch tutorial (0.27\,M) &  51 & 117 & 133 & 318 (4) & 318 & 247 \\
\bottomrule
\end{tabular}}
\caption{Stock v0.37.0, \MC. Every engine's port probes and $S$-parameters were bit-identical to the basic engine's
(patch x5: $S_{11}$ only; notch x3: port~1 only, since the pulse has not reached port~2 by the end of the run).}
\label{tab:step1}
\end{table}

An initial read-only review of the source (kept as \code{results/fable\_optimization\_review.md}) ranked the
candidates by expected gain per effort, using two micro-benchmarks it wrote for the purpose (\code{barrier\_bench.cpp}
and \code{kernel\_bench.cpp}, named there; the later standalone kernels are \code{scripts/kernel/kb\_idx.cpp} and
\code{kb\_tb.cpp}): a barrier diet, the inner loop, dynamic work distribution, the boundary paths, and, further out, a
fused E/H sweep. That order was followed, and the estimates held up within their stated error except that the fused
sweep mattered more and an unfused GPU engine less than predicted. The review's remark that a Metal engine would pay
only with fp16 storage or with temporal blocking was borne out later by the diamond-blocked engine of
Section~\ref{sec:gpu}, which runs the two large models faster than the CPU engine.
